\documentclass[../Main.tex]{subfiles}

\begin{document}
Few quantum mechanical systems can be solved exactly. In perturbation theory we start from a known system (that we can solve),
\begin{equation*}
    H\ket n = E_n \ket n
\end{equation*}
with $\{\ket n\}$ an orthonormal basis of energy eigenstates, and calculate energies and eigenstates for a perturbed Hamiltonian:
\begin{equation}
    (H + \lambda V) \state = E\state
    \label{eqnPerturbedH}
\end{equation}
where $V$ is a Hermitian operator. We assume that $E$ and $\state$ depend smoothly on the parameter $\lambda$, so we can create a power series expansion and taking $\lambda \to 0$ recovers the original system.
\section{Non-Degenerate Case}
Suppose that as $\lambda \to 0$, $\state \to \ket r$ which is a unique eigenstate, and $E \to E_r$ which is a non-degenerate energy level.
We use $\{\ket n\}$ as a basis for $\lambda \neq 0$ so we can write:
\begin{align*}
    \state &= \alpha \ket r + \sum_{j \neq r} \beta_j \ket j \\
    &= \alpha \left(\ket r + \sum_{j \neq r} \gamma_j \ket j\right)
\end{align*}
where $\alpha, \gamma_j = \beta_j / \alpha$ are all power series in $\lambda$ such that $\alpha \to 1$, $\beta_j, \gamma_j \to 0$ when $\lambda \to 0$.

Our aim is to calculate the coefficients in expansions:
\begin{align*}
    E &= E_r + E_r' \lambda + E_r'' \lambda^2 + o(\lambda^2) \\
    \alpha &= 1 + a\lambda + o(\lambda) \\
    \gamma_j &= c_j \lambda + o(\lambda)
\end{align*}
We will first substitute the expression for $\state$ into \eqnref{eqnPerturbedH} and cancel a factor $\alpha$:
\begin{equation*}
    \left(E_r + \lambda V\right) \ket r + \sum_{j \neq r} \gamma_j \left(E_j + \lambda V\right) \ket j = E \left(\ket r + \sum_{j \neq r} \gamma_j \ket j\right)
\end{equation*}
Rearranging:
\begin{equation}
    (E - E_r) \ket r + \sum_{j \neq r} \gamma_j (E - E_j) \ket j = \lambda V \ket r + \lambda \sum_{j \neq r} \gamma_j V \ket j
    \label{eqnSeriesEquation}
\end{equation}
We can now use \eqnref{eqnSeriesEquation} to find equations for the coefficients. Taking $\bra r \times (\text{\ref{eqnSeriesEquation}})$:
\begin{equation*}
    E - E_r = \lambda \opbraket r V r + \lambda \sum_{j \neq r} \gamma_j \opbraket r V j
\end{equation*}
So far this is exact (true for all orders of $\lambda$), so now in terms of series expansion we find:
\begin{equation*}
    \lambda E_r' + \lambda^2 E_r'' + \cdots = \lambda \opbraket r V r + \lambda^2 \sum_{j \neq r}c_j \opbraket r V j + \cdots
\end{equation*}
Thus we see that the first-order change is given by:
\begin{equation*}
    E'_r = \opbraket r V r
\end{equation*}
To find the second-order change, we need to know $c_j$. Considering $\ket i \times (\text{\ref{eqnSeriesEquation}})$:
\begin{equation*}
    \gamma_i (E - E_i) = \lambda\opbraket i V r + \lambda \sum_{j \neq r} \gamma_j \opbraket i V j
\end{equation*}
and again we note this is exact. To find $c_j$, we take the first order term:
\begin{equation*}
    \lambda c_i (E_r - E_i) = \lambda \opbraket i V r
\end{equation*}
which means: $c_i = \frac{\opbraket i V r}{E_r - E_i}$ which very clearly needs non-degeneracy. The second-order change in the energy is:
\begin{equation*}
    E_r'' = \sum_{j\neq r} \frac{\abs{\opbraket j V r}^2}{E_r - E_j}
\end{equation*}
In summary,
\begin{equation}
    E = E_r + \lambda \opbraket r V r + \lambda^2 \sum_{j \neq r} \frac{\abs{\opbraket j V r}^2}{E_r - E_j} + o(\lambda^2)
    \label{eqnNonDegenPerturb}
\end{equation}
\begin{example}
    \label{exHarmOscPerturb}
    We consider a perturbed harmonic oscillator. For this,
    \begin{equation*}
        H = \frac{\hat{p}^2}{2m} + \frac12 m\omega^2 \hat{x}^2 = \hbar \omega \left(a^\dagger a + \frac12\right)
    \end{equation*}
    with states $\ket n$ and energies $E_n = \hbar \omega (n + \frac12)$.

    Let the perturbation be $V = m\omega^2 \hat{x}^2 = \hbar\omega(a^2 + (a^\dagger)^2 + 2a^\dagger a + 1)$. This gives:
    \begin{align*}
        \opbraket n V n &= \frac12 \hbar \omega (2n + 1) \\
        \opbraket {n+2} V n &= \frac12 \hbar \omega \sqrt{(n+1)(n+2)} \\
        \opbraket {n-2} V n &= \frac12 \hbar \omega \sqrt{n(n-1)}
    \end{align*}
    with any other inner product zero.
    
    Now we use \eqnref{eqnNonDegenPerturb}:
    \begin{align*}
        E &\approx E_n + \lambda \opbraket n V n + \lambda^2 \sum_{m \neq n} \frac{\abs{\opbraket m V n}^2}{E_n - E_m} \\
        &= \hbar \omega \left(n + \frac12\right) + \lambda \hbar \omega \left(n + \frac12\right) \\
        &\qquad\qquad+ \lambda^2 \left(\frac{\hbar\omega}{2}\right)^2 \left(\frac{(n+1)(n+2)}{-2\hbar\omega} + \frac{n(n-1)}{2\hbar\omega}\right) \\
        &= \hbar \omega \left(n + \frac12\right) \left(1 + \lambda - \lambda^2 / 2\right)
    \end{align*}
    This is consistent with the transformation $\omega^2 \mapsto \omega^2 (1 + 2\lambda)$ which gives an exact solution for our perturbed problem.
\end{example}
\begin{remarks}
    \item We expect \eqnref{eqnNonDegenPerturb} to be valid and useful when $\abs \lambda$ is sufficiently small. We really have expansion in quantities like:
        \begin{equation*}
            \frac{\lambda^2 \opbraket j V r}{E_i - E_j}
        \end{equation*}
        so accuracy requires the energy shifts (due to perturbation) to be much smaller than the energy differences of the original problem.
    \item The perturbation series may not converge. In this case it can still be useful for approximate calculations as an asymptotic series. In example~\ref{exHarmOscPerturb}, we considered a total potential:
        \begin{equation*}
            V(\hat{x}) = \frac12 m\omega^2 (1 + 2\lambda) \hat{x}^2
        \end{equation*}
        For $\lambda < -\frac12$ we expect the theory not to make sense, which suggests a radius of convergence $\abs{\lambda} < \frac12$.
\end{remarks}
\begin{example}[Ground state for helium]
    The unperturbed problem is two non-interacting electrons in a potential due to a nucleus with charge $2e$:
    \begin{equation*}
        H = H_1(\opvec{x_1}, \opvec{p_1}) + H_2(\opvec{x_1}, \opvec{p_1}),\quad H_1 = H_2 = \frac{\abs{\opvec p}^2}{2m} - \frac{2e^2}{4\pi \epsilon_0 \abs{\opvec x}}.
    \end{equation*}
    One-electron states and energies are given by:
    \begin{equation*}
        \ket{n~l~m},\text{ with energy } E_n = -\frac12 m \left(\frac{2e^2}{4\pi \epsilon_0 \hbar}\right)^2 \frac{1}{n^2}
    \end{equation*}
    The lowest energy eigenstate for two electrons is:
    \begin{equation*}
        \ket\Psi = \ket{1~0~0} \otimes \ket{1~0~0} \otimes \ket\chi
    \end{equation*}
    where $\ket\chi$ is the antisymmetric spin state, $\frac1{\sqrt{2}} \left(\upst \dnst - \dnst\upst\right)$.

    The wavefunction for this two-electron state is:
    \begin{equation*}
        \Psi(\vec{x_1}, \vec{x_2}) = \psi_{100}(\vec{x_1}) \psi_{100}(\vec{x_2})
    \end{equation*}
    where the wavefunction for each electron is:
    \begin{equation*}
        \psi_{100} = \frac{1}{\sqrt{\pi}} \left(\frac{2}{a_2}\right)^{\frac32} e^{-\abs{\vec x} / a_2}
    \end{equation*}
    and the quantity $a_2$ is half the Bohr radius,
    \begin{equation*}
        a_2 = \frac12 \frac{4\pi \epsilon_0}{e^2} \frac{\hbar^2}{m}
    \end{equation*}
    We find the ground state energy, $2E_1$, to be $-108.8~eV$. Experiment shows that the ground state energy for helium is in fact $-79~eV$, so we clearly need the electron-electron interaction. We treat this as a perturbation:
    \begin{equation*}
        \lambda = \frac{e^2}{4\pi \epsilon_0}\quad\text{and}\quad V(\vec{x_1}, \vec{x_2}) = \frac{1}{\abs{\vec{x_1} - \vec{x_2}}}
    \end{equation*}
    The first-order correction to the ground state energy is:
    \begin{align*}
        \lambda \opbraket \Psi V \Psi &= \lambda \int_{\R^3} \int_{\R^3} \Psi(\vec{x_1}, \vec{x_2})^* V(\vec{x_1}, \vec{x_2}) \Psi(\vec{x_1}, \vec{x_2})~d^3\vec{x_2}~d^3\vec{x_2} \\
        &= \lambda \int_{\R^3} \int_{\R^3} \frac{\abs{\psi_{100}(\vec{x_1})}^2\abs{\psi_{100}(\vec{x_2})}^2}{\abs{\vec{x_1} - \vec{x_2}}}~d^3\vec{x_2}~d^3\vec{x_2} \\
        &\approx 34~eV
    \end{align*}
    This gives a corrected ground state energy of $-108.8~eV + 34~eV = -74.8~eV$, which is closer to the true value of $-79~eV$.
\end{example}
\section{Degenerate Case}
Recall the problem:
\begin{equation*}
    (H + \lambda V) \state = E\state
    \tag{\ref{eqnPerturbedH}}
\end{equation*}
Now suppose that as $\lambda \to 0$, $E \to E_r = E_s = \cdots$ with corresponding states $\ket r, \ket s, \cdots$. If we label these states using $r$, This means $\state \to \sum_r a_r \ket r$. Let us label the other basis states $\ket j$, with eigenvalues different from $E_r$.

For general $\lambda$, we expand:
\begin{equation*}
    \psi = \sum_r \alpha_r \ket r + \sum_j \beta_j \ket j
\end{equation*}
and we must have:
\begin{align*}
    \alpha_r &= a_r + O(\lambda) \\
    \beta_r &= O(\lambda)
\end{align*}
Substituting $\state$ into \eqnref{eqnPerturbedH}:
\begin{align*}
    \sum_r \alpha_r (E_r + \lambda V)\ket r &+ \sum_j \beta_j (E_j + \lambda V) \ket j \\
    &= E\left(\sum_r \alpha_r \ker r + \sum_j \beta_j \ket j\right)
\end{align*}
Then rearranging:
\begin{equation}
    \sum_{r} \alpha_r (E - E_r) \ket r + \sum_j \beta_j (E - E_j) \ket j = \lambda \sum_r \alpha_r V \ket{r} + \lambda \sum_j \beta_j V \ket j
    \label{eqnDegenExpand}
\end{equation}
We want $E = E_r + \lambda E' + O(\lambda^2)$, and note that the correction $E'$ is no longer associated with a specific basis state in our chosen state. Taking the inner product of \eqnref{eqnDegenExpand} with $\ket s$ (one of the degenerate states):
\begin{equation*}
    (E - E_s) \alpha_s = \lambda \sum_r \alpha_r \opbraket s v r + \lambda \sum_j \beta_j \opbraket s V j
\end{equation*}
Then to leading order in $\lambda$,
\begin{equation}
    \sum_r \opbraket s V r a_r = E' a_s
    \label{eqnDegenPerturb}
\end{equation}
This is the key result: $E'$ is the eigenvalue of the matrix $\opbraket s V r$ and $a_r$ the components of the eigenvector. If there are $N$ degenerate states, we have $N$ solutions for $E'$ and $a_r$, with corresponding eigenstates $\sum_r a_r \ket r$.

We can interpret these different values for $E'$ as the degenerate energy level $E_r$ splitting apart under the perturbation into new energy levels, breaking degeneracy (or reducing it, if we get repeated solutions for some values of $E'$). Degenerate states usually arise from symmetry of the Hamiltonian, and this perturbation is breaking some/all of this symmetry.
\begin{example}[Particle in a box in 2D]
    Let the position coordinates be $(x, y) \in [0, a]^2$, and our boundary condition is that the wavefunction is zero on the edges. The unperturbed problem has states $\ket{p~q}$ with energies:
    \begin{equation*}
        E_{p, q} = \frac{\hbar^2 \pi^2}{2ma^2} (p^2 + q^2)
    \end{equation*}
    with wavefunctions:
    \begin{equation*}
        \psi_{p, q}(x, y) = \frac2a \sin\left(\frac{p\pi x}{a}\right) \sin \left(\frac{q \pi y}{a}\right)
    \end{equation*}
    The lowest energy level $E_{1, 1}$ is non-degenerate, but the next one up is 2-degenerate: $E_{2, 1} = E_{1, 2}$.

    We perturb as:
    \begin{equation*}
        V(x, y) = \frac{xy}{a^2}
    \end{equation*}
    so the constant $\lambda$ must have dimensions of energy. The shift in the lowest energy is given by \eqnref{eqnNonDegenPerturb} (non-degenerate perturbation):
    \begin{align*}
        \opbraket{1~1}{V}{1~1} &= \int_{0}^{a} \int_{0}^{a} \frac{xy}{a} \abs{\psi_{1, 1}}^2 dy~ dx  \\
        &= \frac14.
    \end{align*}
    The shift in the degenerate energy is given by \eqnref{eqnDegenPerturb} (degenerate perturbation), which requires the eigenvalues of:
    \begin{equation*}
        \begin{pmatrix}
            \opbraket{1~2}{V}{1~2} & \opbraket{1~2}{V}{2~1} \\
            \opbraket{2~1}{V}{1~2} & \opbraket{2~1}{V}{2~1}
        \end{pmatrix}
        =
        \begin{pmatrix}
            \alpha & \beta \\
            \beta & \alpha
        \end{pmatrix}
    \end{equation*}
    and so the eigenvalues of this matrix are $\alpha \pm \beta$ with eigenvectors:
    \begin{equation*}
        \frac{1}{\sqrt{2}} \begin{pmatrix}1 \\ \pm1 \end{pmatrix}
    \end{equation*}
    We calculate $\alpha$ and $\beta$:
    \begin{align*}
        \alpha &= \int_{0}^{a} \int_{0}^{a} \frac{xy}{a^2}\abs{\psi_{12}(x, y)}^2 dy~dx \\
        &= \left(\frac{2}{a}\right)^2 \frac{1}{a^2} \left(\int_{0}^{a} \sin^2 \left(\frac{\pi x}{a}\right) dx\right) \left(\int_{0}^{a} \sin^2 \left(\frac{2\pi y}{a}\right) dy \right) \\
        &= \frac14. \\
        \beta &= \int_{0}^{a} \int_{0}^{a} \frac{xy}{a^2}\psi_{12}(x, y)\psi_{21}(x, y) dy~dx \\
        &= \left(\frac{2}{a}\right)^2 \frac{1}{a^2} \left(\int_{0}^{a} x\sin\left(\frac{\pi x}{a}\right)\sin \left(\frac{2\pi x}{a}\right) dx\right)^2 \\
        &= \left(\frac{4}{3\pi}\right)^4
    \end{align*}
    Summing up, we get the new ground state energy:
    \begin{equation*}
        \frac{\hbar^2 \pi^2}{ma^2} + \frac{\lambda}{4}
    \end{equation*}
    The next two energy levels have energies:
    \begin{equation*}
        \frac{5\hbar^2 \pi^2}{2ma^2} + \frac{\lambda}{4} \pm \lambda \left(\frac{4}{3\pi}\right)^4
    \end{equation*}
    and corresponding states:
    \begin{equation*}
        \frac1{\sqrt{2}} \left(\ket{1~2} \pm \ket{2~1}\right)
    \end{equation*}
\end{example}
\end{document}